\documentclass[conference]{IEEEtran}
\IEEEoverridecommandlockouts
% The preceding line is only needed to identify funding in the first footnote. If that is unneeded, please comment it out.
\usepackage{cite}
\usepackage{amsmath,amssymb,amsfonts}
\usepackage{algorithmic}
\usepackage{graphicx}
\usepackage{textcomp}
\usepackage{xcolor}
\usepackage{hyperref}
\def\BibTeX{{\rm B\kern-.05em{\sc i\kern-.025em b}\kern-.08em
    T\kern-.1667em\lower.7ex\hbox{E}\kern-.125emX}}
\begin{document}

\title{Laboratório 2 - Unidade Operativa Uniciclo RISC-V32I\\}

\author{\IEEEauthorblockN{1\textsuperscript{st} Artur Henrique do Carmo Rosa}
\IEEEauthorblockA{\textit{Departamento de Ciências da Computação} \\
\textit{Universidade de Brasília}\\
Brasília, Brazil \\
222028738@aluno.unb.br}
\and
\IEEEauthorblockN{2\textsuperscript{nd} Daniel Dezan Baby}
\IEEEauthorblockA{\textit{Departamento de Ciências da Computação} \\
\textit{Universidade de Brasília}\\
Brasília, Brazil \\
232046097@aluno.unb.br}
\and
\IEEEauthorblockN{3\textsuperscript{rd} João Filipe Araújo}
\IEEEauthorblockA{\textit{Departamento de Ciências da Computação} \\
\textit{Universidade de Brasília}\\
Brasília, Brazil \\
231013574@aluno.unb.br}
\and
\IEEEauthorblockN{4\textsuperscript{th} Leonardo Tomé Sampaio}
\IEEEauthorblockA{\textit{Departamento de Ciências da Computação} \\
\textit{Universidade de Brasília}\\
Brasília, Brazil \\
222005448@aluno.unb.br}
\and
\IEEEauthorblockN{5\textsuperscript{th} Luís Eduardo Curi Serra}
\IEEEauthorblockA{\textit{Departamento de Ciências da Computação} \\
\textit{Universidade de Brasília}\\
Brasília, Brazil \\
251033574@aluno.unb.br}
}

\maketitle

\begin{abstract}
No presente trabalho é apresentado o desenvolvimento da Unidade Operativa Uniciclo da arquitetura
RISC-V32I, implementada no ambiente Quartus II. O projeto contempla os módulos que compõem o caminho de
dados (\textit{datapath}), as memórias de instrução e de dados em organização Harvard, a unidade de
controle principal e a unidade de controle da ULA, de forma a executar automaticamente um conjunto de
instruções de 32 bits pré-compiladas no formato MIF. Além do conjunto básico de referência
(Tipo-R, \texttt{lw}, \texttt{sw} e \texttt{beq}), foram implementadas as instruções de salto e de
formação de constantes (\texttt{jal}, \texttt{jalr}, \texttt{lui} e \texttt{auipc}), o que exigiu a
extensão do caminho de dados clássico com novos sinais de controle e novos multiplexadores. O trabalho
inclui a descrição da arquitetura desenvolvida, do conjunto de instruções suportado, dos procedimentos e
métodos adotados, dos resultados obtidos por meio do ambiente de simulação Waveform Editor e da análise
de desempenho das unidades funcionais, seguidos de uma análise crítica dos resultados e das principais
dificuldades de implementação e depuração enfrentadas, em especial nas instruções de salto e no
tratamento de endereços fora da região válida da memória de instruções.
\end{abstract}

\begin{IEEEkeywords}
RISC-V, Uniciclo, FPGA, Quartus, caminho de dados, unidade de controle, multiplexador, saltos, jal, jalr
\end{IEEEkeywords}

%-----------------------------------------------------------------------

\section{Divisão de tarefas}

\subsection{Objetivos e Introdução (Teoria da Arquitetura)}
Seção ii. Objetivos: resumo sucinto dos objetivos do laboratório.
Seção iii. Introdução (parte 1): visão geral da arquitetura RISC-V32I, filosofia RISC, organização
Uniciclo e caminho de dados.

\subsection{Introdução (Caminho de Dados e Memórias)}
Seção iii. Introdução (parte 2): detalhamento dos módulos do caminho de dados, memórias de instrução e
dados em organização Harvard, banco de registradores, somador, gerador de imediato e ULA.

\subsection{Introdução (Unidade de Controle e Multiplexadores)}
Seção iii. Introdução (parte 3): unidade de controle principal, unidade de controle da ULA, sinais de
controle e multiplexadores do caminho de dados.

\subsection{Procedimentos e Métodos}
Seção iv. Materiais e Métodos: ferramentas utilizadas (Quartus II, Waveform Editor), descrição da
implementação de cada módulo e metodologia de sincronização e integração.

\subsection{Resultados experimentais}
Seção v. Resultados: simulações no Waveform Editor, conteúdo dos registradores, sinais de controle e
análise de desempenho de cada unidade funcional.

\subsection{Análise de resultados, Conclusões e Revisão Geral}
Seção vi. Discussão e Conclusões: análise qualitativa e quantitativa, gargalos identificados, pontos
críticos, dificuldades de implementação e depuração, e conclusão geral do laboratório.

%-----------------------------------------------------------------------

\section{Objetivos}
O objetivo deste laboratório foi desenvolver a Unidade Operativa Uniciclo da arquitetura RISC-V32I,
utilizando o ambiente Quartus II Altera, de forma a montar os módulos principais do caminho de dados
(\textit{datapath}) capazes de processar instruções de 32 bits pré-compiladas no formato MIF, fazendo uso
de memórias de instrução e de dados em organização Harvard. Além disso, o trabalho buscou reforçar os
conceitos estudados na disciplina de Organização e Arquitetura de Computadores, proporcionando uma
compreensão mais prática sobre o funcionamento do caminho de dados, da unidade de controle e da análise de
desempenho de um processador Uniciclo. De forma específica, buscou-se estender o modelo de referência
apresentado em sala \cite{patterson2020computer} para dar suporte às instruções de salto incondicional
(\texttt{jal}, \texttt{jalr}) e às instruções de formação de constantes (\texttt{lui}, \texttt{auipc}),
exercitando a inclusão de novos sinais de controle e novos caminhos de dados na microarquitetura.

\section{Introdução}

\subsection{Arquitetura RISC-V32I Uniciclo}
A arquitetura RISC-V32I segue a filosofia RISC (\textit{Reduced Instruction Set Computer}), caracterizada
por um conjunto reduzido e regular de instruções de tamanho fixo (32 bits), poucos formatos de codificação
e um modelo \textit{load--store}, no qual apenas as instruções de carga e armazenamento acessam a memória
de dados, enquanto as demais operam exclusivamente sobre o banco de registradores
\cite{patterson2020computer, referenceGuide}. A regularidade desse conjunto facilita a decodificação e o
projeto do caminho de dados, pois os campos da instrução (opcode, \textit{funct3}, \textit{funct7},
registradores e imediato) ocupam posições previsíveis em cada formato.

Na organização Uniciclo, toda instrução é integralmente buscada, decodificada e executada dentro de um
único ciclo de \textit{clock}. Não há sobreposição entre estágios, de modo que o período do \textit{clock}
deve ser dimensionado para acomodar o pior caminho crítico — tipicamente o caminho da instrução de carga
(\texttt{lw}), que percorre memória de instruções, banco de registradores, ULA e memória de dados antes da
escrita de volta \cite{patterson2020computer}. Como consequência, a simplicidade de controle do modelo
Uniciclo é obtida ao custo de um período de \textit{clock} ditado pela instrução mais lenta.

O projeto adota organização de memória do tipo Harvard, com memória de instruções e memória de dados
fisicamente separadas, cada uma com seu próprio barramento de endereços e de dados. Essa separação permite
que a busca da instrução e o eventual acesso ao dado ocorram no mesmo ciclo, sem conflito de acesso a uma
única memória. O caminho de dados é composto pelo contador de programa (PC), pela memória de instruções,
pelo banco de registradores, pelo gerador de imediato, pela Unidade Lógica e Aritmética (ULA) com sua
unidade de controle, pelos somadores de cálculo de endereço, pela memória de dados e por um conjunto de
multiplexadores que roteiam os dados conforme os sinais produzidos pela unidade de controle principal.

\subsection{Contador de Programa e Memória de Instruções}
O contador de programa (PC) é um registrador de 32 bits que armazena o endereço da instrução em execução.
Trata-se do único elemento de estado do estágio de busca: a cada borda de subida do \textit{clock}, o PC é
atualizado com o endereço da próxima instrução, calculado de forma combinacional pelo restante do caminho
de dados. O módulo possui ainda uma entrada de \textit{reset} assíncrono que, quando ativada, inicializa o
PC com o endereço base do segmento de texto, \texttt{0x00004000}, alinhado ao mapa de memória do RISC-V32I
adotado pelo simulador RARS, conforme recomendado na especificação do trabalho.

A memória de instruções é descrita como uma memória somente de leitura no fluxo de execução, inicializada a
partir do arquivo \texttt{UnicicloInst.mif}. Como o PC trabalha com endereços de \textit{byte} a partir de
\texttt{0x00004000}, e a memória é organizada em palavras de 32 bits, é necessário traduzir o endereço
recebido em um índice de palavra: subtrai-se o endereço base (\texttt{TEXT\_BASE}) do endereço corrente e
divide-se o resultado por quatro (deslocamento de dois bits à direita). Além disso, a memória verifica se o
endereço solicitado pertence à faixa válida do segmento de texto; caso contrário, retorna uma instrução
sentinela definida, evitando que a leitura de uma posição não inicializada propague valores indefinidos
para o restante do caminho de dados. Esse mecanismo de validação é discutido em detalhe na
Seção~\ref{sec:dificuldades}.

\subsection{Banco de Registradores}
O banco de registradores implementa os 32 registradores de 32 bits do RISC-V32I (\texttt{x0} a
\texttt{x31}). Ele dispõe de duas portas de leitura assíncronas (combinacionais), endereçadas pelos campos
\textit{rs1} e \textit{rs2} da instrução, e de uma porta de escrita síncrona, endereçada pelo campo
\textit{rd} e habilitada pelo sinal de controle \texttt{RegWrite}. A escrita ocorre na borda de subida do
\textit{clock}, garantindo que o valor lido em um ciclo corresponda ao conteúdo anterior à escrita do mesmo
ciclo. O registrador \texttt{x0} é mantido permanentemente em zero: sua leitura sempre retorna zero e
qualquer tentativa de escrita nele é ignorada, conforme a especificação da arquitetura
\cite{referenceGuide}. Para fins de verificação, o módulo expõe ainda o conteúdo de todos os 32
registradores em saídas dedicadas, possibilitando a visualização simultânea de toda a janela de
registradores na simulação, atendendo à exigência da especificação de apresentar em tela o conteúdo dos
trinta e dois registradores.

\subsection{Unidade Lógica e Aritmética (ULA) e seu Controle}
A Unidade Lógica e Aritmética (ULA) é o bloco combinacional responsável pelas operações aritméticas e
lógicas do processador. Em sua entrada recebe dois operandos de 32 bits e um código de operação de 4 bits;
em sua saída produz o resultado de 32 bits e o sinal \texttt{Zero}, que indica se o resultado é nulo e é
utilizado pela lógica de desvio condicional. As operações implementadas incluem soma, subtração, AND, OR,
XOR, deslocamento lógico à esquerda (\texttt{SLL}), deslocamento lógico à direita (\texttt{SRL}),
comparação com sinal (\texttt{SLT}) e uma operação de repasse do segundo operando (\texttt{PASSA-B}),
introduzida especialmente para dar suporte à instrução \texttt{lui}.

A escolha da operação efetiva fica a cargo da unidade de controle da ULA (\textit{ALU control}), um bloco
combinacional que combina o sinal \texttt{ALUOp}, gerado pela unidade de controle principal, com os campos
\textit{funct3} e o bit~30 da instrução (que distingue, por exemplo, \texttt{add} de \texttt{sub}). Esse
arranjo em dois níveis evita que a unidade de controle principal precise conhecer todos os detalhes de cada
operação: ela apenas indica a categoria (acesso à memória, desvio, Tipo-R, Tipo-I ou \texttt{lui}), e a
unidade de controle da ULA refina essa categoria na operação final \cite{patterson2020computer}.

\subsection{Somador}
O caminho de dados utiliza somadores dedicados de 32 bits, descritos como blocos puramente combinacionais
que produzem a soma de duas entradas. São empregados dois somadores no estágio de atualização do PC: o
primeiro calcula \texttt{PC + 4}, endereço da instrução sequencial seguinte; o segundo calcula
\texttt{PC + imediato}, endereço alvo dos desvios condicionais (\texttt{beq}, \texttt{bne}) e do salto
\texttt{jal}. Manter esses somadores como módulos separados e simples mantém o caminho de dados próximo do
modelo de referência e facilita o rastreamento dos endereços na simulação.

\subsection{Gerador de Imediato}
O gerador de imediato (\textit{Imm Gen}) extrai e monta o valor imediato de 32 bits a partir dos campos da
instrução, selecionando o formato correto com base no opcode. A arquitetura RISC-V32I codifica os imediatos
em diferentes formatos (I, S, B, U e J), nos quais os bits do imediato aparecem espalhados em posições
distintas da palavra de instrução, sempre de modo a manter o bit de sinal na posição mais significativa
\cite{referenceGuide}. O módulo realiza a extensão de sinal para 32 bits em todos os formatos com sinal e,
nos formatos de desvio (B) e de salto (J), reintroduz o bit menos significativo igual a zero, refletindo o
fato de que esses deslocamentos são múltiplos de dois \textit{bytes}. A correta montagem desses 32 bits é
um ponto sensível do projeto e está diretamente associada às dificuldades relatadas na
Seção~\ref{sec:dificuldades}, uma vez que um imediato mal formado leva ao cálculo de endereços de desvio
incorretos.

\subsection{Memória de Dados}
A memória de dados armazena os operandos manipulados pelas instruções de carga (\texttt{lw}, \texttt{lhu})
e armazenamento (\texttt{sw}), sendo inicializada a partir do arquivo \texttt{UnicicloData.mif}. Ela é
endereçada pelo resultado da ULA (que, nessas instruções, corresponde ao endereço efetivo \texttt{rs1 +
imediato}) e controlada pelos sinais \texttt{MemRead} e \texttt{MemWrite}. Para compatibilizar a leitura e
a escrita dentro de um único ciclo de \textit{clock}, adotou-se a estratégia de escrever na memória na
borda de subida e disponibilizar o dado lido na borda de descida do \textit{clock}, de modo que a leitura e
a escrita não disputem o mesmo instante de tempo e o caminho de dados permaneça coerente em uma execução
Uniciclo.

\subsection{Unidade de Controle Principal}
A Unidade de Controle Uniciclo é o bloco responsável por observar o campo opcode da instrução, com 7 bits
de largura, e decidir quais sinais devem ser ativados no restante do caminho de dados para que aquela
instrução específica seja executada corretamente \cite{patterson2020computer}. Diferente de um registrador
ou de uma memória, ela não guarda nenhum valor de um ciclo para o outro: assim que a memória de instruções
entrega um novo opcode na sua entrada, a saída já reflete os novos sinais de controle, respeitando apenas
o tempo de propagação da lógica combinacional utilizada na sua implementação.

Os sinais produzidos por esse bloco, junto com o efeito de cada um sobre o caminho de dados, são os
seguintes:

\begin{itemize}
    \item RegWrite: habilita a escrita no banco de registradores;
    \item MemRead: habilita a leitura da memória de dados;
    \item MemWrite: habilita a escrita na memória de dados;
    \item MemtoReg: seleciona, por meio de multiplexadores, a origem do valor gravado no registrador de
    destino; foi ampliado para 2 bits a fim de incluir, além do resultado da ULA e do dado lido da memória,
    o endereço de retorno \texttt{PC + 4} usado por \texttt{jal} e \texttt{jalr};
    \item ALUSrc: seleciona, por meio de um multiplexador, se a segunda entrada da ULA é o conteúdo
    do segundo registrador lido ou o valor imediato gerado pelo gerador de imediato;
    \item Branch: indica que a instrução é um desvio condicional, sinal que é combinado com a saída Zero
    da ULA (e com o campo \textit{funct3}) para decidir o próximo valor do contador de programa;
    \item ALUOp: indica de forma resumida o tipo de operação que a ULA deve realizar, sendo depois
    combinado com os campos funct3 e funct7 pela unidade de controle da ULA para definir a operação final;
    \item Jump, JumpR e PCtoALU: sinais adicionais, ausentes no modelo de referência, introduzidos para dar
    suporte às instruções \texttt{jal}, \texttt{jalr} e \texttt{auipc}, respectivamente, conforme detalhado
    na Seção~\ref{sec:dificuldades}.
\end{itemize}

A Tabela \ref{tab:controle} apresenta a tabela verdade reduzida da unidade de controle para o modelo de
referência, relacionando cada grupo de instruções, identificado por seu opcode, aos valores dos sinais de
controle gerados \cite{patterson2020computer}. O valor X (don't care) indica situações em que o sinal não
influencia o resultado da instrução, pois seu efeito não é utilizado naquele caso. No modelo clássico, o
campo ALUOp utiliza 2 bits para identificar a categoria da operação a ser realizada pela ULA: 00 para
acessos à memória, 01 para desvios e 10 para operações lógico-aritméticas. Na implementação deste trabalho,
o campo ALUOp foi ampliado para 3 bits, a fim de também distinguir as operações com imediato (Tipo-I) e a
operação de repasse exigida pela instrução \texttt{lui}, sendo o detalhamento da operação final feito pela
unidade de controle da ULA a partir dos campos funct3 e funct7. A tabela completa, com todas as instruções
implementadas e os sinais adicionais, é apresentada na Seção de Procedimentos
(Tabela~\ref{tab:controle_completo}).

\begin{table}[!t]
\caption{Tabela verdade reduzida da Unidade de Controle (modelo de referência)}
\begin{center}
\scriptsize
\setlength{\tabcolsep}{3pt}
\begin{tabular}{c|c|c|c|c|c|c|c|c}
\hline \hline
\textbf{Instr.} & \textbf{Opcode} & \textbf{ASrc} & \textbf{M2Reg} & \textbf{RegW} & \textbf{MemR} & \textbf{MemW} & \textbf{Br} & \textbf{ALUOp} \\
\hline
Tipo-R & 0110011 & 0 & 0 & 1 & 0 & 0 & 0 & 10 \\
lw & 0000011 & 1 & 1 & 1 & 1 & 0 & 0 & 00 \\
sw & 0100011 & 1 & X & 0 & 0 & 1 & 0 & 00 \\
beq & 1100011 & 0 & X & 0 & 0 & 0 & 1 & 01 \\
\hline \hline
\end{tabular}
\end{center}
\label{tab:controle}
\end{table}

Esses quatro grupos representam o modelo de referência apresentado em sala \cite{patterson2020computer}. A
mesma estrutura combinacional foi utilizada para reconhecer também os demais opcodes do conjunto de
instruções adotado no projeto, atribuindo a cada um deles a combinação adequada de sinais conforme o
formato correspondente. Para opcodes não previstos, todos os sinais de habilitação são mantidos desligados
por padrão, evitando que uma instrução inválida cause efeitos indesejados no restante do caminho de dados.

\subsection{Multiplexador 2:1}
O multiplexador 2 para 1 é o elemento combinacional mais simples utilizado no caminho de dados: ele recebe
duas entradas de dados e, dependendo do valor de um único bit de seleção, decide qual delas aparece na
saída \cite{patterson2020computer}. A Tabela \ref{tab:mux} resume seu comportamento.

\begin{table}[!t]
\caption{Tabela verdade do Multiplexador 2:1}
\begin{center}
\begin{tabular}{c|c}
\hline \hline
\textbf{Seleção} & \textbf{Saída} \\
\hline
0 & Entrada 0 \\
1 & Entrada 1 \\
\hline \hline
\end{tabular}
\end{center}
\label{tab:mux}
\end{table}

Em vez de montar um circuito seletor diferente para cada ponto do caminho de dados onde havia essa
necessidade, optou-se por um único módulo genérico e parametrizável quanto à largura do barramento, depois
instanciado nos diversos pontos do caminho de dados Uniciclo em que existe seleção entre duas fontes de
dados. O primeiro deles escolhe entre o dado lido do segundo registrador e o valor imediato, de acordo com
o sinal ALUSrc. Outro multiplexador escolhe entre o conteúdo de \textit{rs1} e o valor do PC como primeiro
operando da ULA, em apoio à instrução \texttt{auipc}. No estágio de escrita de volta, dois
multiplexadores em cascata selecionam entre o resultado da ULA, o dado lido da memória e o endereço de
retorno \texttt{PC + 4}, de acordo com o sinal MemtoReg de 2 bits. Por fim, no estágio de atualização do
PC, multiplexadores em cascata escolhem entre o endereço sequencial \texttt{PC + 4}, o endereço de desvio
\texttt{PC + imediato} e o endereço de \texttt{jalr}, conforme a combinação dos sinais Branch/Zero,
\texttt{Jump} e \texttt{JumpR}.

\section{Procedimentos e Métodos}

\subsection{Ferramentas Utilizadas}
O desenvolvimento foi realizado no ambiente Quartus II Altera, utilizando descrição de \textit{hardware} em
linguagem Verilog para todos os módulos do caminho de dados. A verificação funcional foi conduzida no
módulo de simulação Waveform Editor (Vector Waveform Editor) do próprio Quartus, complementada, em alguns
módulos, por \textit{testbenches} escritos em Verilog que comparam automaticamente a saída obtida com o
valor esperado em cada caso de teste. As memórias de instruções e de dados foram inicializadas a partir dos
arquivos no formato MIF (\texttt{UnicicloInst.mif} e \texttt{UnicicloData.mif}), de acordo com o mapa de
endereços do RISC-V32I sugerido na especificação. A metodologia adotada foi incremental: cada módulo foi
descrito e validado isoladamente antes de ser integrado ao módulo de topo (\texttt{riscv\_cpu}), o que
reduziu o espaço de busca durante a depuração da unidade operativa completa.

\subsection{Contador de Programa e Memória de Instruções}
O contador de programa foi descrito como um registrador de 32 bits sensível à borda de subida do
\textit{clock}, com \textit{reset} assíncrono que o inicializa em \texttt{0x00004000}. A cada ciclo, o PC
recebe o endereço calculado pela lógica de seleção do próximo PC. A memória de instruções foi descrita como
um vetor de palavras de 32 bits inicializado pelo arquivo MIF. A leitura é combinacional: a partir do
endereço de \textit{byte} fornecido pelo PC, calcula-se o índice de palavra subtraindo o endereço base e
deslocando dois bits à direita. Implementou-se ainda uma verificação de faixa (\textit{valid\_addr}): caso
o endereço esteja abaixo do início do segmento de texto ou além do tamanho da memória, a saída é forçada
para a instrução sentinela \texttt{0x0000006F} (equivalente a \texttt{jal x0, 0}), que mantém o PC preso na
posição corrente em vez de buscar uma palavra não inicializada. A motivação dessa decisão de projeto é
discutida na Seção~\ref{sec:dificuldades}.

\subsection{Banco de Registradores}
O banco foi descrito como um vetor de 32 posições de 32 bits, com leitura combinacional nas portas
\textit{rs1} e \textit{rs2} e escrita síncrona na borda de subida do \textit{clock}, condicionada ao sinal
\texttt{RegWrite} e ao fato de o registrador de destino não ser \texttt{x0}. Tanto na leitura quanto na
escrita, \texttt{x0} é tratado como constante igual a zero. Para atender ao requisito de visualização, o
conteúdo de cada registrador foi também conectado a saídas individuais nomeadas conforme a convenção ABI
do RISC-V (\texttt{ra}, \texttt{sp}, \texttt{gp}, \texttt{tp}, \texttt{t0}--\texttt{t6},
\texttt{s0}--\texttt{s11}, \texttt{a0}--\texttt{a7}), permitindo acompanhar a janela completa de
registradores durante a simulação.

\subsection{Unidade Lógica e Aritmética e seu Controle}
A ULA foi descrita como um bloco combinacional cujo comportamento é selecionado por um código de 4 bits.
Além das operações aritmético-lógicas usuais (soma, subtração, AND, OR, XOR, deslocamentos e comparação com
sinal), foi incluída uma operação de repasse do segundo operando (\texttt{PASSA-B}), utilizada pela
instrução \texttt{lui}. A unidade de controle da ULA foi descrita de forma hierárquica: a partir do sinal
\texttt{ALUOp} (3 bits), seleciona-se a categoria da operação; quando necessário, os campos \textit{funct3}
e o bit~30 da instrução refinam a operação final. Em particular, distinguiu-se o caso das instruções com
imediato (Tipo-I), nas quais \texttt{funct3 = 000} deve sempre resultar em soma (\texttt{addi}), do caso
Tipo-R, no qual o bit~30 diferencia \texttt{add} de \texttt{sub}. Ambos os módulos foram validados
isoladamente, aplicando-se todas as combinações relevantes de \texttt{ALUOp}, \textit{funct3} e bit~30 e
conferindo o código de operação e o resultado produzidos.

\subsection{Somador}
Os somadores foram descritos como blocos combinacionais de 32 bits e instanciados duas vezes no caminho de
dados, para o cálculo de \texttt{PC + 4} e de \texttt{PC + imediato}. A validação consistiu em verificar,
para diferentes valores de PC e de imediato, se a saída correspondia exatamente à soma esperada,
incluindo casos com imediato negativo (desvios para trás), nos quais a extensão de sinal do imediato é
essencial para o cálculo correto do endereço.

\subsection{Gerador de Imediato}
O gerador de imediato foi descrito como um bloco combinacional que seleciona o formato (I, S, B, U ou J)
a partir do opcode e monta o valor de 32 bits por concatenação dos campos correspondentes, com extensão de
sinal. Tomou-se cuidado especial para que cada concatenação resultasse em exatamente 32 bits e para que os
formatos de desvio (B) e de salto (J) reintroduzissem o bit menos significativo igual a zero. A verificação
foi feita aplicando-se instruções representativas de cada formato e conferindo o imediato resultante, com
atenção particular aos casos de sinal negativo e aos deslocamentos de \texttt{beq}/\texttt{bne} e
\texttt{jal}.

\subsection{Memória de Dados}
A memória de dados foi descrita como um vetor de palavras de 32 bits inicializado pelo arquivo MIF,
endereçado pelo resultado da ULA. A escrita ocorre na borda de subida do \textit{clock}, condicionada ao
sinal \texttt{MemWrite}, e a leitura é capturada na borda de descida, condicionada implicitamente ao uso do
dado pela instrução. Essa separação temporal entre escrita e leitura dentro do mesmo ciclo foi adotada para
manter a coerência da execução Uniciclo, evitando condições de corrida entre as duas operações.

\subsection{Unidade de Controle Principal}
A unidade de controle principal foi descrita como um bloco de lógica combinacional, utilizando uma
estrutura de decisão sobre o campo opcode de 7 bits extraído diretamente da instrução lida da memória.
Para cada opcode suportado, foi definido um conjunto fixo de valores para os sinais de saída (RegWrite,
MemRead, MemWrite, MemtoReg, ALUSrc, Branch, ALUOp, Jump, JumpR e PCtoALU), seguindo a tabela verdade
completa apresentada na Tabela~\ref{tab:controle_completo}. Para opcodes não previstos, todos os sinais de
habilitação foram mantidos desligados por padrão, evitando que uma instrução inválida cause algum efeito
indesejado no restante do caminho de dados.

A validação desse módulo foi feita de forma isolada no Waveform Editor do Quartus, aplicando-se cada
opcode como entrada e conferindo, um a um, se a combinação de sinais gerada na saída correspondia
exatamente à tabela de projeto antes da integração com o restante do caminho de dados.

\begin{table*}[!t]
\caption{Tabela verdade completa da Unidade de Controle implementada}
\begin{center}
\scriptsize
\setlength{\tabcolsep}{4pt}
\begin{tabular}{l|c|c|c|c|c|c|c|c|c|c|c}
\hline \hline
\textbf{Instr.} & \textbf{Opcode} & \textbf{ALUSrc} & \textbf{Mem2Reg} & \textbf{RegWrite} & \textbf{MemRead} & \textbf{MemWrite} & \textbf{Branch} & \textbf{Jump} & \textbf{JumpR} & \textbf{PCtoALU} & \textbf{ALUOp} \\
\hline
Tipo-R                  & 0110011 & 0 & 00 & 1 & 0 & 0 & 0 & 0 & 0 & 0 & 010 \\
Tipo-I (addi, andi, \ldots) & 0010011 & 1 & 00 & 1 & 0 & 0 & 0 & 0 & 0 & 0 & 011 \\
lw / lhu                & 0000011 & 1 & 01 & 1 & 1 & 0 & 0 & 0 & 0 & 0 & 000 \\
sw                      & 0100011 & 1 & XX & 0 & 0 & 1 & 0 & 0 & 0 & 0 & 000 \\
beq / bne               & 1100011 & 0 & XX & 0 & 0 & 0 & 1 & 0 & 0 & 0 & 001 \\
jal                     & 1101111 & X & 10 & 1 & 0 & 0 & 0 & 1 & 0 & 0 & 000 \\
jalr                    & 1100111 & 1 & 10 & 1 & 0 & 0 & 0 & 0 & 1 & 0 & 000 \\
lui                     & 0110111 & 1 & 00 & 1 & 0 & 0 & 0 & 0 & 0 & 0 & 100 \\
auipc                   & 0010111 & 1 & 00 & 1 & 0 & 0 & 0 & 0 & 0 & 1 & 000 \\
\hline
inválido (default)      & ------- & 0 & 00 & 0 & 0 & 0 & 0 & 0 & 0 & 0 & 000 \\
\hline \hline
\end{tabular}
\end{center}
\label{tab:controle_completo}
\end{table*}

\subsection{Multiplexador 2:1}
O módulo do multiplexador foi escrito de forma genérica, com a largura do barramento de dados como
parâmetro, o que permitiu reaproveitá-lo nos diversos pontos de seleção do caminho de dados sem reescrever a
lógica em cada lugar. Essa escolha reduz a duplicação de código e facilita a verificação, já que o
comportamento do multiplexador só precisa ser validado uma vez para que a confiança se estenda a todas as
instâncias seguintes. A validação foi feita no Waveform Editor, alternando o sinal de seleção entre 0 e 1
para diferentes combinações de entrada e conferindo se a entrada correta era propagada para a saída em
cada caso.

\subsection{Integração e Lógica de Salto}
A integração dos módulos foi feita no módulo de topo, conectando os sinais de controle às entradas de
seleção dos multiplexadores. Para suportar as instruções de salto e de formação de constantes, três pontos
do caminho de dados foram acrescentados em relação ao modelo de referência: (i) um multiplexador que
seleciona o primeiro operando da ULA entre \textit{rs1} e o PC, controlado por \texttt{PCtoALU} e usado pela
instrução \texttt{auipc}; (ii) um segundo nível no multiplexador de escrita de volta, que permite gravar
\texttt{PC + 4} no registrador de destino para \texttt{jal} e \texttt{jalr}; e (iii) uma cascata de
multiplexadores na atualização do PC, que seleciona entre \texttt{PC + 4}, o endereço de desvio
\texttt{PC + imediato} (compartilhado por \texttt{beq}, \texttt{bne} e \texttt{jal}) e o endereço de
\texttt{jalr}, este último obtido de \texttt{(rs1 + imediato)} com o bit menos significativo forçado a zero.
A decisão de desvio condicional foi isolada em um módulo dedicado (\textit{BranchControl}), que combina o
sinal \texttt{Branch}, a saída \texttt{Zero} da ULA e o campo \textit{funct3} para distinguir \texttt{beq}
(salta se igual) de \texttt{bne} (salta se diferente).

\section{Resultados experimentais}

\subsection{Contador de Programa e Memória de Instruções}
Na simulação, verificou-se que, após o \textit{reset}, o PC parte de \texttt{0x00004000} e avança de quatro
em quatro \textit{bytes} a cada ciclo, exceto quando uma instrução de desvio ou salto altera o fluxo. A
memória de instruções entregou, a cada endereço, a palavra de 32 bits correspondente do arquivo
\texttt{UnicicloInst.mif}, e a verificação de faixa comportou-se conforme o esperado: endereços fora do
segmento de texto resultaram na instrução sentinela \texttt{0x0000006F}, mantendo o PC estável em vez de
buscar conteúdo indefinido.

\subsection{Banco de Registradores}
A simulação confirmou a leitura assíncrona nas portas \textit{rs1} e \textit{rs2} e a escrita síncrona em
\textit{rd}, com \texttt{x0} permanecendo em zero em todas as situações. A exposição individual dos 32
registradores permitiu acompanhar, lado a lado, o efeito de cada instrução sobre o estado da máquina,
facilitando a conferência dos resultados das demais unidades.

\subsection{Unidade Lógica e Aritmética e seu Controle}
A ULA e sua unidade de controle produziram, para cada combinação de \texttt{ALUOp}, \textit{funct3} e
bit~30, o código de operação e o resultado esperados, incluindo o caso especial \texttt{PASSA-B} usado pela
instrução \texttt{lui} e a distinção entre \texttt{addi} (soma forçada) e o par \texttt{add}/\texttt{sub}.
O sinal \texttt{Zero} acompanhou corretamente o resultado nulo, dando suporte à decisão de desvio
condicional.

\subsection{Somador}
Os somadores produziram \texttt{PC + 4} e \texttt{PC + imediato} corretamente, inclusive para imediatos
negativos, situação em que o desvio para trás depende da extensão de sinal correta do imediato.

\subsection{Gerador de Imediato}
O gerador de imediato reproduziu os valores esperados para os formatos I, S, B, U e J, com a extensão de
sinal correta e com o bit menos significativo igual a zero nos formatos B e J. A conferência desses valores
foi determinante para validar os endereços de desvio e de salto, conforme discutido na
Seção~\ref{sec:dificuldades}.

\subsection{Memória de Dados}
A memória de dados respondeu corretamente às instruções \texttt{lw} e \texttt{sw}: a escrita na borda de
subida e a leitura na borda de descida permitiram que, dentro de um mesmo ciclo, o dado armazenado fosse
posteriormente recuperado sem condição de corrida, mantendo a coerência da execução Uniciclo.

\subsection{Unidade de Controle Principal e Multiplexador 2:1}
A verificação de ambos os módulos foi feita por meio de \textit{testbenches} próprios, escritos em Verilog
e executados no simulador, que comparam automaticamente a saída obtida com o valor esperado em cada caso de
teste e reportam o resultado.

Para a unidade de controle, o \textit{testbench} aplicou, um a um, os opcodes correspondentes a todos os
grupos de instruções suportados pelo projeto, além de um opcode inválido para verificar o comportamento
padrão. A Tabela \ref{tab:result_controle} resume os casos testados, todos aprovados.

\begin{table}[!t]
\caption{Casos de teste da Unidade de Controle}
\begin{center}
\begin{tabular}{c|c|c}
\hline \hline
\textbf{Opcode (bin)} & \textbf{Opcode (hex)} & \textbf{Resultado} \\
\hline
0110011 & 33 & PASS \\
0010011 & 13 & PASS \\
0000011 & 03 & PASS \\
0100011 & 23 & PASS \\
1100011 & 63 & PASS \\
1101111 & 6F & PASS \\
1100111 & 67 & PASS \\
0110111 & 37 & PASS \\
0010111 & 17 & PASS \\
1111111 & 7F & PASS \\
\hline \hline
\end{tabular}
\end{center}
\label{tab:result_controle}
\end{table}

Como a unidade de controle é um circuito puramente combinacional, a forma de onda gerada na simulação
apresenta os sinais de saída constantes para cada opcode aplicado, sem transições intermediárias, o que é
o comportamento esperado para esse tipo de lógica.

Para o multiplexador, o \textit{testbench} aplicou diferentes combinações de dados de 32 bits nas duas
entradas, variando o sinal de seleção entre 0 e 1, e conferiu se a saída reproduzia exatamente a entrada
selecionada em cada caso. Foram usados valores representativos como AAAAAAAA, 55555555, DEADBEEF e
CAFEBABE, e todos os casos foram aprovados, confirmando o funcionamento correto do módulo para qualquer
dado de entrada.

% TODO (grupo): inserir capturas de tela do Waveform Editor com as formas de onda da unidade de controle,
% do multiplexador e da execução integrada (PC, instrução, registradores e sinais de controle), conforme
% exigido no Requisito 4 da especificação.

\subsection{Análise de Desempenho das Unidades Funcionais}
% TODO (grupo): preencher com os tempos medidos no modo de compilação Timing do Quartus.
A análise de desempenho deve ser obtida no modo de compilação \textit{Timing} do Quartus, que considera os
atrasos de propagação reais de cada unidade funcional. Qualitativamente, espera-se que o caminho crítico do
processador Uniciclo seja determinado pela instrução de carga (\texttt{lw}), cujo percurso atravessa, em
sequência, a memória de instruções, o banco de registradores, o gerador de imediato, a ULA e a memória de
dados, antes da escrita de volta no banco de registradores. As instruções de salto (\texttt{jal},
\texttt{jalr}) e de formação de constantes (\texttt{lui}, \texttt{auipc}) tendem a apresentar caminhos mais
curtos, por não acessarem a memória de dados. A Tabela de tempos por módulo (valor absoluto em
submúltiplos de segundo e percentual do ciclo) e a estimativa do período de \textit{clock} devem ser
preenchidas a partir dos relatórios de temporização da síntese, identificando os pontos de
\textit{gargalo} (\textit{bottlenecks}).

\section{Análise de Resultados}
A unidade de controle principal e o multiplexador atenderam ao comportamento esperado em todos os casos de
teste, gerando os sinais corretos para cada opcode e propagando o dado certo em função do sinal de seleção.
A verificação automática por \textit{testbench} foi importante porque permitiu confirmar o comportamento de
cada módulo de forma isolada, antes da integração com o restante do caminho de dados, reduzindo a chance de
erros difíceis de localizar em uma etapa posterior. Após a integração, o caminho de dados completo executou
o conjunto de instruções pré-compiladas, atualizando corretamente o PC, o banco de registradores e a
memória de dados ao longo dos ciclos.

\subsection{Dificuldades de Implementação e Depuração}
\label{sec:dificuldades}
As principais dificuldades do trabalho concentraram-se em dois eixos relacionados: (i) a implementação das
instruções de salto e de formação de constantes (\texttt{jal}, \texttt{jalr}, \texttt{lui} e
\texttt{auipc}), que extrapolam o caminho de dados de referência; e (ii) o tratamento de saltos para fora
da região válida da memória de instruções, agravado por imediatos que inicialmente não eram montados
corretamente em 32 bits.

\subsubsection{Implementação dos saltos, \texttt{auipc}, \texttt{lui}, \texttt{jal} e \texttt{jalr}}
O caminho de dados Uniciclo de referência \cite{patterson2020computer} contempla apenas as instruções
Tipo-R, \texttt{lw}, \texttt{sw} e \texttt{beq}. A inclusão das demais instruções exigiu acrescentar
sinais de controle e caminhos de dados que não existem na Figura~1 da especificação, o que foi a fonte da
maior parte das dificuldades:

\begin{itemize}
    \item \textbf{Endereço de retorno (\texttt{jal}/\texttt{jalr}).} Essas instruções precisam gravar
    \texttt{PC + 4} no registrador de destino, mas o multiplexador de escrita de volta original só escolhia
    entre o resultado da ULA e o dado da memória. Foi necessário ampliar o sinal \texttt{MemtoReg} para 2
    bits e acrescentar um segundo nível de multiplexador para rotear \texttt{PC + 4} até o banco de
    registradores. Antes dessa correção, o registrador de destino recebia o resultado da ULA, e os
    endereços de retorno gravados estavam errados, o que só se manifestava ciclos depois, ao retornar do
    salto.

    \item \textbf{\texttt{auipc} — PC como operando da ULA.} A instrução \texttt{auipc} calcula
    \texttt{PC + imediato}, e não \texttt{rs1 + imediato}. Como a ULA, no modelo de referência, recebe
    sempre \textit{rs1} em seu primeiro operando, foi preciso introduzir o sinal \texttt{PCtoALU} e um
    multiplexador adicional que substitui \textit{rs1} pelo PC nessa instrução.

    \item \textbf{\texttt{lui} — repasse do imediato.} A instrução \texttt{lui} deveria simplesmente
    escrever o imediato (com os 12 bits inferiores zerados) no registrador de destino. Uma armadilha sutil
    é que os bits \texttt{[19:15]} da instrução \texttt{lui} fazem parte do imediato, mas o caminho de dados
    sempre os interpreta como o endereço de \textit{rs1}; assim, somar \textit{rs1} ao imediato produziria
    um resultado contaminado por uma leitura ``fantasma'' de registrador. A solução foi criar uma operação
    dedicada na ULA (\texttt{PASSA-B}, código \texttt{4'b1000}), selecionada por \texttt{ALUOp = 100}, que
    descarta o primeiro operando e repassa apenas o imediato. O mesmo cuidado vale para \texttt{auipc}, em
    que a leitura fantasma de \textit{rs1} é neutralizada pelo multiplexador de \texttt{PCtoALU}.

    \item \textbf{Cascata de seleção do próximo PC.} Com três origens possíveis para o próximo endereço
    (\texttt{PC + 4}; \texttt{PC + imediato} para \texttt{beq}/\texttt{bne}/\texttt{jal}; e o alvo de
    \texttt{jalr}), foi necessário organizar uma cascata de multiplexadores com a prioridade correta. O
    primeiro nível escolhe entre \texttt{PC + 4} e \texttt{PC + imediato} a partir de
    \texttt{(PcSrc OR Jump)}; o segundo nível sobrepõe o alvo de \texttt{jalr} quando \texttt{JumpR} está
    ativo. Erros na ordem dessa cascata faziam, por exemplo, um \texttt{jalr} ser ``mascarado'' por um
    desvio condicional inativo.

    \item \textbf{\texttt{jalr} — alinhamento do alvo.} O endereço de destino do \texttt{jalr} é
    \texttt{(rs1 + imediato)} com o bit menos significativo forçado a zero
    (\texttt{\{alu\_result[31:1], 1'b0\}}). Omitir esse zeramento gerava endereços desalinhados, que por
    sua vez caíam fora das posições válidas da memória de instruções.

    \item \textbf{Distinção \texttt{jal} $\times$ \texttt{jalr} no gerador de imediato.} Embora ambas sejam
    saltos incondicionais, \texttt{jal} usa imediato no formato J (relativo ao PC) e \texttt{jalr} usa
    imediato no formato I (relativo a \textit{rs1}, com \texttt{ALUSrc = 1}). Tratar as duas com o mesmo
    formato de imediato levava a alvos completamente incorretos.
\end{itemize}

\subsubsection{Saltos para fora da memória e validação dos 32 bits}
A segunda grande dificuldade surgiu justamente quando os saltos passaram a funcionar: em algumas execuções,
o PC alcançava endereços fora da região efetivamente carregada do segmento de texto. Como, inicialmente, a
memória de instruções não validava o endereço recebido, a leitura de uma posição não inicializada produzia
uma palavra de 32 bits indefinida (valores \texttt{X} na simulação). Essa palavra indefinida era então
decodificada pela unidade de controle como um opcode arbitrário, ativando sinais de controle aleatórios e
corrompendo o banco de registradores e a memória de dados. O sintoma observado — registradores assumindo
valores sem sentido vários ciclos após o salto — era difícil de localizar, pois a causa (o salto para fora
da memória) e o efeito (a corrupção do estado) estavam separados no tempo.

A investigação revelou duas causas que se reforçavam mutuamente:

\begin{enumerate}
    \item \textbf{Tradução de endereço ausente.} Como o PC inicia em \texttt{0x00004000} e não em zero, o
    índice de palavra não pode ser obtido diretamente de \texttt{addr >> 2}; é preciso primeiro subtrair o
    endereço base do segmento (\texttt{byte\_offset = addr - TEXT\_BASE}) e só então deslocar dois bits à
    direita. Sem essa tradução, qualquer endereço já indexava posições muito além do conteúdo carregado.

    \item \textbf{Imediatos não validados em 32 bits.} O código de montagem dos imediatos, em versões
    iniciais, não garantia que cada concatenação somasse exatamente 32 bits — seja por erro na largura da
    replicação do bit de sinal, seja por omissão do bit menos significativo implícito (igual a zero) nos
    formatos B e J. Um imediato mal formado produzia um \textit{offset} de desvio/salto incorreto e, por
    consequência, fazia o desvio aterrissar fora da região correta da memória. Esse é o elo direto entre o
    problema da ``validação dos 32 bits'' e os ``saltos para fora da memória''.
\end{enumerate}

A solução combinou três medidas. Primeiro, padronizou-se a tradução de endereço relativa a
\texttt{TEXT\_BASE}. Segundo, reescreveu-se o gerador de imediato verificando, campo a campo, que cada
formato soma exatamente 32 bits: por exemplo, o formato B é montado como
$20 + 1 + 6 + 4 + 1 = 32$ bits e o formato J como $12 + 8 + 1 + 10 + 1 = 32$ bits, sempre com o bit menos
significativo igual a zero. Terceiro, adicionou-se à memória de instruções uma verificação de faixa
(\textit{valid\_addr}): quando o endereço está abaixo do início do segmento de texto ou além do tamanho da
memória, a saída é forçada para a instrução sentinela \texttt{0x0000006F} (\texttt{jal x0, 0}), que prende
o PC na posição corrente. Essa instrução sentinela transforma um salto para fora da memória — antes
catastrófico e silencioso — em uma condição controlada e observável: em vez de propagar valores
indefinidos, a máquina entra em um laço estável, tornando o problema imediatamente visível na simulação.

A metodologia de depuração que permitiu isolar essas falhas baseou-se na verificação modular (cada bloco
testado isoladamente antes da integração) e na observação, ciclo a ciclo, do PC e da palavra de instrução
de 32 bits no Waveform Editor. A introdução da instrução sentinela foi particularmente útil porque deu uma
``assinatura'' reconhecível ao evento de salto para fora da memória, encurtando o tempo entre a manifestação
do sintoma e a identificação da causa.

\section{Conclusão}
A construção da Unidade Operativa Uniciclo RISC-V32I mostrou, na prática, como a decodificação do opcode
pode ser isolada em um bloco puramente combinacional, como um único componente seletor genérico pode ser
reaproveitado em vários pontos do caminho de dados e como a separação física entre memórias de instrução e
de dados (organização Harvard) viabiliza a busca e o acesso a dados em um único ciclo. Essa organização
modular deixou o projeto mais claro e facilitou a verificação individual de cada parte antes da integração
da unidade operativa completa.

As maiores dificuldades não estiveram nos módulos do modelo de referência, mas na extensão necessária para
suportar as instruções de salto e de formação de constantes (\texttt{jal}, \texttt{jalr}, \texttt{lui} e
\texttt{auipc}), que demandaram novos sinais de controle (\texttt{Jump}, \texttt{JumpR}, \texttt{PCtoALU},
\texttt{MemtoReg} de 2 bits), uma nova operação na ULA (\texttt{PASSA-B}) e uma cascata de multiplexadores
para a seleção do próximo PC. O caso mais instrutivo foi o dos saltos para fora da memória de instruções:
ele evidenciou como um imediato mal validado em 32 bits e a ausência de tradução e validação de endereço
podem se combinar em uma falha de difícil diagnóstico, e como uma decisão de projeto simples — retornar uma
instrução sentinela bem definida para endereços inválidos — pode converter uma falha silenciosa em um
comportamento estável e observável. O laboratório reforçou, assim, tanto os conceitos de organização e
arquitetura de computadores quanto a importância de uma metodologia de verificação incremental e de
decisões de projeto defensivas.

\section*{Agradecimentos}



\begin{thebibliography}{00}
\bibitem{referenceGuide}
F. Vidal,
\textit{RISC-V Reference Guide v25}, 2025.

\bibitem{patterson2020computer}
D.~A. Patterson e J.~L. Hennessy,
\textit{Computer Organization and Design RISC-V Edition: The Hardware Software Interface}.
Morgan Kaufmann, 2020.

\end{thebibliography}

\end{document}
